% 第 16 题原料用量表：按固定源 PDF 矢量路径换算（1 pt → 0.03515 cm）；原点取表格左下（源页 (913.89, 512.81) pt），y 向上。
% 4 条横线（含上下外框）、5 条竖线；表头行高 30.63 pt 约为数据行 15.51 pt 的两倍；
% 角格 45° 斜线自左上至右下，右上「原料」（对应 A,B,C 三列）、左下「肥料」（对应甲、乙两行）。
% 列宽（自左至右）66.28 / 20.22 / 20.84 / 22.85 pt；列标 A,B,C 为数学斜体，数据 4/8/3/5/5/10 为正体数字。
\begin{tikzpicture}[line cap=round,line join=round,font=\small,x=0.03515cm,y=0.03515cm]
% 横线（4 条）
\draw (0,61.65)--(130.59,61.65);
\draw (0,31.02)--(130.59,31.02);
\draw (0,15.51)--(130.59,15.51);
\draw (0,0)--(130.59,0);
% 竖线（5 条）
\draw (0.20,0)--(0.20,61.65);
\draw (66.48,0)--(66.48,61.65);
\draw (86.70,0)--(86.70,61.65);
\draw (107.54,0)--(107.54,61.65);
\draw (130.39,0)--(130.39,61.65);
% 角格斜线（左上→右下）
\draw (0.40,61.45)--(66.29,31.22);
% 斜线表头两角文字
\node[anchor=base west] at (39.32,50.87) {原料};
\node[anchor=base west] at (6.38,35.76) {肥料};
% 表头列标与数据行（各列居中）
\node[anchor=base] at (76.59,43.71) {\(A\)};
\node[anchor=base] at (97.12,43.71) {\(B\)};
\node[anchor=base] at (118.97,43.71) {\(C\)};
\node[anchor=base] at (33.34,20.24) {甲};
\node[anchor=base] at (76.59,20.24) {\(4\)};
\node[anchor=base] at (97.12,20.24) {\(8\)};
\node[anchor=base] at (118.97,20.24) {\(3\)};
\node[anchor=base] at (33.34,4.73) {乙};
\node[anchor=base] at (76.59,4.73) {\(5\)};
\node[anchor=base] at (97.12,4.73) {\(5\)};
\node[anchor=base] at (118.97,4.73) {\(10\)};
\end{tikzpicture}
